\documentclass[a4paper,11pt]{article}
\pdfoutput=1

\usepackage[table]{xcolor}
\usepackage{lmodern}
\usepackage{jcappub}
\usepackage[toc,page]{appendix}
\usepackage{subcaption}
\usepackage{mathtools}
\usepackage{orcidlink}
\usepackage{multirow}
\usepackage[T1]{fontenc}
\usepackage{float}
\usepackage{bm}
\usepackage{tcolorbox}

\tcbuselibrary{skins}
\tcbuselibrary{breakable}

\newcommand{\semibold}[1]{{\fontseries{b}\selectfont{#1}}}
\newcommand{\sansbold}[1]{{\sffamily\fontseries{sbc}\selectfont{#1}}}

\newcommand{\para}[1]{\par\vspace{2mm}\noindent\semibold{{#1.}---}\ignorespaces}

\DeclareMathOperator{\Or}{O}
\DeclareMathOperator{\tr}{tr}
\DeclareMathOperator{\NewRe}{Re}
\DeclareMathOperator{\NewIm}{Im}
\renewcommand{\Re}{\NewRe}
\renewcommand{\Im}{\NewIm}

\newcommand{\transpose}{\mathsf{T}}

\renewcommand{\leq}{\leqslant}
\renewcommand{\geq}{\geqslant}

\definecolor{SussexCobaltBlue}{HTML}{1d4289}
\definecolor{SussexDeepAquamarine}{HTML}{007a78}
\definecolor{SussexPowderBlue}{HTML}{7da1c4}
\definecolor{SussexCornYellow}{HTML}{f2c75c}
\definecolor{SussexChinaRose}{HTML}{be84a3}
\definecolor{SussexBurntOrange}{HTML}{dc582a}

\hypersetup{colorlinks=true,citecolor=SussexDeepAquamarine,linkcolor=SussexCobaltBlue,urlcolor=SussexBurntOrange}

\newcommand{\etal}{\emph{et al.}}

\renewcommand{\d}{\mathrm{d}}
\newcommand{\e}[1]{\mathrm{e}^{{#1}}}
\newcommand{\im}{\mathrm{i}}
\newcommand{\Mp}{M_{\mathrm{P}}}
\newcommand{\vect}[1]{\bm{\mathrm{{#1}}}}

\DeclareMathOperator{\grad}{\nabla}
\DeclareMathOperator{\sinc}{sinc}

\newcommand{\GeV}{\text{GeV}}

% spread out arrays / text
\renewcommand{\arraystretch}{1.15}
\renewcommand{\baselinestretch}{1.1}

\newenvironment{panel}{\begin{tcolorbox}[enhanced,breakable,colback=gray!10,colframe=gray!30]\small}{\end{tcolorbox}}
\newenvironment{todo}{\begin{tcolorbox}[enhanced,breakable,colback=SussexBurntOrange!10,colframe=SussexBurntOrange!40]\small{\textsf{\textbf{Open issue:} }}}{\end{tcolorbox}}
\newenvironment{calculation}{\begin{tcolorbox}[enhanced,breakable,colback=SussexPowderBlue!10,colframe=SussexCobaltBlue!30]\small}{\end{tcolorbox}}

%
% ----------------------------------------------------------------------
% Notation macros. Collected here so that definitions can be changed
% in one place. Grouped by role.
% ----------------------------------------------------------------------
%

% -- e-folding time and background quantities --
\newcommand{\Ny}{N}                          % e-fold "time" variable
\newcommand{\Ninit}{N_{\mathrm{init}}}       % start of instanton transition
\newcommand{\Nfinal}{N_{\mathrm{final}}}     % end of instanton transition
\newcommand{\DN}{\Delta N}                   % duration of transition, N_init - N_final on background
\newcommand{\dNstar}{\delta N_\star}         % instanton excess e-folds
\newcommand{\epsSR}{\epsilon_1}              % first slow-roll parameter
\newcommand{\Hub}{H}                         % Hubble rate
\newcommand{\sfactor}{a}                     % scale factor
\newcommand{\atloc}{(\field,\mom)}           % shorthand: evaluated at the local phase-space
                                              % point (field,mom); distinguishes the LOCAL,
                                              % dynamical use of H, eps_1 (Friedmann-equation
                                              % functions of the local trajectory, exactly as
                                              % in AbstractPotential.H_sq/epsilon) from the
                                              % GLOBAL background aH(N) used only to define
                                              % the coordinate y and the gradient-term prefactor

% -- field content --
\newcommand{\field}{\phi}                    % inflaton field
\newcommand{\mom}{\pi}                       % conjugate momentum, d(field)/dN
\newcommand{\Vpot}{V}                        % potential
\newcommand{\phinl}{\field_{\mathrm{nl}}}    % noiseless background trajectory
\newcommand{\pinl}{\mom_{\mathrm{nl}}}       % noiseless background momentum
\newcommand{\phiinit}{\field_{\mathrm{init}}}
\newcommand{\phiend}{\field_{\mathrm{end}}}

% -- noise / MSR response fields --
\newcommand{\noise}[1]{\xi_{#1}}             % generic noise, e.g. \noise{\field}
\newcommand{\rfield}{\widetilde\field}       % MSR response field conjugate to \field
\newcommand{\rmom}{\widetilde\mom}           % MSR response field conjugate to \mom

% -- shell / layer bookkeeping --
\newcommand{\Nlay}[1]{N_{#1}}                % comoving population of shell i, N_i
\newcommand{\rad}{r}                         % comoving radial coordinate

% -- the onion coordinate y and its complement u --
\newcommand{\yco}{y}                         % onion coordinate, y=-1 fixed outer edge,
                                              % y=+1 current horizon (exterior-only, log-radial)
\newcommand{\uco}{u}                         % u = 1-y (legacy; unused in current coordinate)
\newcommand{\rout}{\rad_{\mathrm{out}}}      % comoving radius of outer edge of overdensity
\newcommand{\rphys}{\rad_{\mathrm{phys}}}    % present-day physical scale (Leach-Liddle), distinct
                                              % from comoving \rad and areal R (Section~\ref{sec:scale-assignment})
\newcommand{\aHo}{(\sfactor \Hub)_0}         % aH evaluated at N_init
\newcommand{\aH}[1][]{(\sfactor \Hub)#1}     % aH, optionally with an argument e.g. \aH{(N)}
\newcommand{\aHloc}{\sfactor(\Ny)\Hub\atloc} % LOCAL comoving-Hubble-rate quantity for the
                                              % gradient (Laplacian) coupling term: a(N) is a
                                              % pure function of the grid N (N is itself each
                                              % shell's own e-fold clock, so a(N)=a_0 e^{N-Ninit}
                                              % is exact, not tied to any trajectory), while H
                                              % is local, H(field,mom). Contrast \aH[(\Ny)],
                                              % retained for the coordinate/geometric roles
                                              % (defining y, aHo, n(y,N), k_sigma) where a single
                                              % shared comoving map is required.
\newcommand{\rH}{\rad_H}                     % current horizon radius, evaluated at the CORE's
                                              % own local field values: rH(N) = 1/[a(N)H(1,N)]
\newcommand{\Deltas}{\Delta s}                % Deltas(N) = ln(rout/rH(N)): log-radial extent
                                              % of the resolved (exterior) domain
\newcommand{\epscore}{\epsSR^{\mathrm{core}}} % epsilon_1 evaluated at the core, y=+1
\newcommand{\measure}{\mu}                   % self-adjoint measure for Lop(y,N): exponential
                                              % in y, N-dependent through Deltas(N)
\newcommand{\advcoef}{A}                     % advection coefficient in (y,N) coordinates,
                                              % nonzero for the log-radial horizon-excising map
\newcommand{\weightfn}{w}                    % w(y) = mu(y,N)*p(y), where p=e^{Deltas y} is the
                                              % coefficient of Q'' in Lop; equivalently mu = w^3.
                                              % Used in the self-adjointness check (Section~\ref{sec:msr-action})

% -- LEGACY: sinc-eigenfunction spectral expansion, superseded by the
% collocation scheme (Section~\ref{sec:collocation-scheme}). Retained,
% unused, only because a handful of macro names below are reused for
% collocation-grid quantities with a different meaning; do not rely on
% the "mode coefficient" comments below, which describe the old scheme.
\newcommand{\Lop}{\mathcal{L}_\yco}          % radial mean-field Laplacian in y (now the
                                              % log-radial, horizon-excising operator of
                                              % Section~\ref{sec:coordinate}, not the old sinc one)
\newcommand{\modeidx}{n}                     % mode index
\newcommand{\basisfn}{f_\modeidx}            % sinc basis function
\newcommand{\basisfnx}[1]{f_{#1}}            % sinc basis function, explicit index
\newcommand{\eigval}{\kappa_\modeidx}        % eigenvalue of \Lop, = -(n pi)^2
\newcommand{\modenorm}{N_\modeidx}           % norm of basis function f_n in the (1-y)^2 dy measure

% -- mode coefficients --
\newcommand{\coefa}{a_\modeidx}              % mode coefficient of (field - noiseless)
\newcommand{\coefax}[1]{a_{#1}}              % mode coefficient of (field - noiseless), explicit index
\newcommand{\coefb}{b_\modeidx}              % mode coefficient of (mom - noiseless)
\newcommand{\rcoefa}{\widetilde{a}_\modeidx} % mode coefficient of response field \rfield
\newcommand{\rcoefax}[1]{\widetilde{a}_{#1}} % mode coefficient of \rfield, explicit index
\newcommand{\rcoefb}{\widetilde{b}_\modeidx} % mode coefficient of response field \rmom

\newcommand{\rcoefbx}[1]{\widetilde{b}_{#1}} % mode coefficient of \rmom, explicit index
\newcommand{\projone}{[1]_\modeidx}          % projection of the constant function 1 onto
                                              % mode n in the u^2 du measure; = 2(-1)^{n+1}
% -- noise amplitude / covariance functions --
\newcommand{\Dnoise}[1]{D_{#1}}              % D_field(y,N), D_mom(y,N), D_{field mom}(y,N):
                                              % shell-averaged noise variance/covariance
                                              % densities (the "2D" convention: matches the
                                              % factor of 2 multiplying D11,D12,D22 in the
                                              % existing FullInstanton forward-pass equations)
\newcommand{\Dcm}[1]{\mathsf{D}_{#1}}        % single-patch (undiluted) diffusion matrix
                                              % element in the codebase convention, i.e.
                                              % DiffusionModel.D_matrix(phi,pi) -> (D11,D12,D22);
                                              % \Pspec{fg} = 2\,\Dcm{fg}\atloc
\newcommand{\Pspec}[1]{P_{#1}}               % reduced power spectra of Briaud et al.
\newcommand{\ncount}{n}                      % Hubble-volume count per unit shell, n(y,N)

% -- misc --
\newcommand{\lag}{\lambda}                   % Lagrange multiplier for core endpoint condition
\newcommand{\zprofile}{\zeta}                % curvature perturbation profile

\title{The onion model: a mean-field gradient-corrected instanton\\
for the spacetime embedding of PBH-forming trajectories}

\author{David Seery}
\affiliation{Astronomy Centre, University of Sussex, Brighton BN1 9QH, United Kingdom}
\emailAdd{D.Seery@sussex.ac.uk}

\begin{document}
\maketitle

% \begin{abstract}
% \noindent
% These notes record the derivation of a mean-field, gradient-corrected
% extension of the instanton method used elsewhere in this project to compute
% primordial black hole (PBH) formation probabilities in stochastic inflation.
% In the existing prescription (`\texttt{FullInstanton}'/`\texttt{SlowRollInstanton}'),
% a single instanton trajectory is assumed to describe the assembly of the
% central region of a spherically symmetric overdensity, while surrounding
% Hubble volumes are assumed to detach and evolve noiselessly once they exit
% the active core. This neglects the gradient coupling between the detaching
% shell and the central region identified by Briaud \etal\ in their treatment
% of gradient corrections to stochastic inflation. Here we build a mean-field
% `onion model': angle-averaged Langevin equations for the field and momentum
% in concentric comoving shells, including the leading gradient coupling
% between shells, converted to a Martin--Siggia--Rose (MSR) representation,
% with the gradient-corrected instanton found as the saddle point of a global
% boundary value problem. We derive the appropriate comoving coordinate for
% the shell structure, construct the MSR action with the correct volume
% measure, derive the instanton (Euler--Lagrange) equations and their
% boundary conditions, construct an analytic eigenbasis for the spatial
% operator, and set out a spectral collocation scheme for the numerical
% solution. Two open theoretical issues are flagged explicitly and are not
% yet resolved: the correct probabilistic interpretation of the two-dimensional
% MSR action relative to the existing single-trajectory instanton, and the
% precise sense in which the terminal (`largest time') boundary condition on
% the MSR response fields can be placed at an arbitrary time later than the
% end of the transition.
% \end{abstract}

\tableofcontents

% ======================================================================
\section{Motivation}
\label{sec:motivation}
% ======================================================================

In the existing instanton approach to PBH formation, a single trajectory
$\field(\Ny)$, $\mom(\Ny) = \d\field/\d\Ny$ is found as the saddle point of
the Martin--Siggia--Rose (MSR) action for the central Hubble patch, subject
to boundary conditions fixing the initial field value $\phiinit$, the final
field value $\phiend$, and the duration $\dNstar$ of the anomalous
(noise-driven) excursion relative to the noiseless background. The
resulting curvature perturbation profile $\zprofile(\rad)$ is built by
assuming that, at each step of the assembly, the current Hubble volume
splits into a central region --- which continues to be described by the
instanton --- and a surrounding shell, which detaches and thereafter
evolves \emph{noiselessly} along the background trajectory.

This is a reasonable first approximation, but it neglects the gradient
coupling between the detaching shell and the central region. Briaud
\etal\ have shown that adjacent Hubble volumes do not decouple completely
upon separation: the noise driving neighbouring patches is correlated on
the coarse-graining scale, and this correlation sources a residual gradient
interaction that persists for of order one $e$-fold after separation. This
coupling makes it more costly, in the sense of MSR action, to build a
sharply peaked (`Cerro Torre') profile than the noiseless-detachment
prescription would suggest, because the detaching shell resists being left
behind by the central assembly.

\begin{panel}
The Briaud \etal\ formalism itself is not directly applicable to the
rare-event geometry we are interested in, because it is built for
statistically typical patches relaxing towards the background, whereas here
the central patch is conditioned to be an extreme fluctuation. The
construction below is therefore not simply an import of their results:
we use their identification of the leading gradient coupling and noise
correlation structure, but the conditioning is handled correctly by
extremizing a single joint action for the whole configuration, rather than
by substituting an unconditional noise correlation function into a Langevin
equation by hand.
\end{panel}

The goal of the `onion model' developed here is to replace the ad hoc
noiseless-detachment rule with a dynamical boundary condition derived from
an action principle, so that the cost of assembling a sharp profile ---
and the associated smoothing of the profile by gradient coupling to the
surrounding shells --- emerges from the saddle point rather than being
imposed by hand.

% ======================================================================
\section{Single-patch Langevin system and shell averaging}
\label{sec:shell-averaging}
% ======================================================================

\subsection{The un-coarse-grained system}

Working to the order in the gradient expansion retained by Briaud \etal,
the Langevin equations for the coarse-grained field at a comoving point
$\vect{x}$ are
\begin{subequations}
\label{eq:patch-langevin}
\begin{align}
    \frac{\d \field}{\d \Ny}(\vect{x}, \Ny)
    & =
    \mom(\vect{x}, \Ny) + \noise{\field}(\vect{x}, \Ny) ,
    \\
    \frac{\d \mom}{\d \Ny}(\vect{x}, \Ny)
    & =
    -(3-\epsSR) \mom(\vect{x},\Ny)
    - \frac{\Vpot'(\field)}{\Hub^2}
    + \frac{1}{\aHloc^2} \grad^2 \field(\vect{x}, \Ny)
    + \noise{\mom}(\vect{x}, \Ny) ,
\end{align}
\end{subequations}
where $\grad^2$ is the comoving spatial Laplacian and the noises have
covariance
\begin{equation}
    \big\langle \noise{f}(\vect{x}, \Ny) \, \noise{g}(\vect{y}, \Ny') \big\rangle
    =
    \Pspec{fg}(\Ny) \, \delta_D(\Ny - \Ny') \, (1-\epsSR) \,
    \sinc\!\big( k_\sigma |\vect{x}-\vect{y}| \big) ,
    \qquad f, g \in \{\field, \mom\} ,
    \label{eq:noise-covariance}
\end{equation}
with $k_\sigma = \sigma \aH$ the coarse-graining scale ($\sigma \ll 1$) and
$\Pspec{fg}$ the reduced power spectra. The noise is white in $\Ny$ and
correlated between comoving points only within the coarse-graining radius
$k_\sigma^{-1}$.

\subsection{Shell averaging}
\begin{panel}
\semibold{$H$, $\epsSR$, $\Pspec{fg}$ here are local, not background,
quantities -- and so is $H$ inside the gradient-term coefficient.}
Throughout \eqref{eq:patch-langevin}--\eqref{eq:noise-covariance}
$\Hub$, $\epsSR$, and $\Pspec{fg}$ should be read as the usual
Friedmann-equation functions of the local phase-space point
$\field(\vect{x},\Ny),\mom(\vect{x},\Ny)$ -- exactly as in the single-patch
stochastic-$\delta N$ formalism this section is built on, and as
\texttt{AbstractPotential.H\_sq}/\texttt{epsilon} and
\texttt{DiffusionModel.D\_matrix} already treat them in the codebase --
\emph{not} as fixed background values. This is unremarkable at the
single-patch level, where there is only one field value in play, but it
matters once shells with genuinely different field content are compared
below and, later, once the shell/mode structure is built in
$(\yco,\Ny)$: see the explicit local-evaluation prescription in
Section~\ref{sec:noise-y} and the corrected equations of
Section~\ref{sec:instanton-eqs}, which this un-coarse-grained system should
be understood to feed into consistently. The same applies to the $\aH$
appearing in the gradient term $\aHloc^{-2}\grad^2\field$: $a(\Ny)$ is a
pure function of the grid $\Ny$ (exact, not an approximation, since $\Ny$
is by construction each patch's own e-fold clock -- $a(\Ny)=a_0\e{\Ny-\Ninit}$
regardless of trajectory), but $H$ there is the same local
$H\atloc$, not a value pinned to any reference trajectory. This is
distinct from the genuinely background $\aHo,\aH[(\Ny)]$ used later purely
to define the comoving coordinate $\yco$
(Section~\ref{sec:coordinate}), which necessarily stay
tied to a single global choice since they set up one shared comoving
labelling for the whole configuration -- see the discussion following
\eqref{eq:gradient-term-y}.
\end{panel}


Define the shell-averaged field for the set of comoving patches making up
shell $i$ as
\begin{equation}
    \field_i(\Ny) \equiv \frac{1}{\Nlay{i}} \sum_{\vect{x} \in i} \field(\vect{x}, \Ny) ,
\end{equation}
where $\Nlay{i}$ is the number of coarse-graining patches contained in
shell $i$. Averaging \eqref{eq:patch-langevin} over the shell, and
dropping the difference between $\langle \Vpot'(\field)\rangle$ and
$\Vpot'(\field_i)$ as a higher-order mean-field correction (suppressed by
$\Vpot''' \times \mathrm{Var}[\field]/\Nlay{i}$, itself already
$\Or(1/\Nlay{i})$ -- and, by the same argument, similarly approximating
$H^2,\epsSR$ evaluated on the shell by their value at the shell-mean
$\field_i,\mom_i$), gives
\begin{subequations}
\begin{align}
    \frac{\d \field_i}{\d \Ny}
    & =
    \mom_i + \overline{\noise{\field}}^{(i)} ,
    \\
    \frac{\d \mom_i}{\d \Ny}
    & =
    -\big(3-\epsSR(\field_i,\mom_i)\big)\mom_i - \frac{\Vpot'(\field_i)}{\Hub^2(\field_i,\mom_i)}
    + \frac{1}{\sfactor(\Ny)^2\Hub^2(\field_i,\mom_i)}\big\langle \grad^2\field \big\rangle_i
    + \overline{\noise{\mom}}^{(i)} .
\end{align}
\end{subequations}

\subsection{Central limit suppression of the shell-averaged noise}

The variance of the shell-averaged noise is
\begin{equation}
    \mathrm{Var}\big[\overline{\noise{\field}}^{(i)}\big]
    =
    \frac{1}{\Nlay{i}^2} \sum_{\vect{x},\vect{y} \in i}
    \Pspec{\field\field} \, \delta_D(\Ny-\Ny') \, \sinc(k_\sigma|\vect{x}-\vect{y}|) .
\end{equation}
If shell $i$ has linear comoving extent much larger than the
coarse-graining radius $k_\sigma^{-1}$, the double sum is dominated by
$\Or(\Nlay{i})$ correlated pairs out of $\Nlay{i}^2$ total, giving the
central-limit-theorem result
\begin{equation}
    \mathrm{Var}\big[\overline{\noise{\field}}^{(i)}\big]
    \; \approx \;
    \frac{c \, \Pspec{\field\field}}{\Nlay{i}} \, \delta_D(\Ny-\Ny') ,
    \label{eq:CLT-suppression}
\end{equation}
for some $\Or(1)$ geometric constant $c$, and similarly for
$\overline{\noise{\mom}}^{(i)}$.

\subsection{Inter-shell correlation is nearest-neighbour and subleading}

The covariance between the averaged noise in shells $i$ and $j$ is
\begin{equation}
    \mathrm{Cov}\big[\overline{\noise{\field}}^{(i)}, \overline{\noise{\field}}^{(j)}\big]
    =
    \frac{1}{\Nlay{i}\Nlay{j}} \sum_{\vect{x}\in i, \vect{y} \in j}
    \Pspec{\field\field} \, \delta_D(\Ny-\Ny') \, \sinc(k_\sigma|\vect{x}-\vect{y}|) .
\end{equation}
Since the sinc kernel connects only points within the coarse-graining
radius $k_\sigma^{-1}$, and the shell-to-shell separation is set by the
comoving Hubble radius --- parametrically larger than $k_\sigma^{-1}$
because $\sigma \ll 1$ --- this covariance vanishes for non-adjacent shells,
and for adjacent shells is suppressed both by the CLT factor
$1/\sqrt{\Nlay{i}\Nlay{j}}$ and by the fraction, $\Or(\Nlay{i}^{-1/3})$, of a
shell's patches that lie within the coarse-graining radius of its boundary.

\begin{panel}
Counting more carefully: the correlated part of the double sum has
$\Or(\Nlay{i}^{2/3})$ contributing pairs (those within a coarse-graining
length of the shared boundary) rather than $\Or(\Nlay{i}^2)$. Relative to
the leading CLT term \eqref{eq:CLT-suppression}, the correlation correction
within a single shell is therefore suppressed by an additional power
$\Nlay{i}^{-1/3}$. The inter-shell noise correlation is thus a genuinely
subleading effect once shells are averaged; the surviving inter-shell
coupling that matters at leading order is the \emph{deterministic} gradient
(Laplacian) term, discussed next.
\end{panel}

\subsection{The mean-field gradient (Laplacian) coupling}

The shell-averaged Laplacian, for a spherically decomposed field, is the
standard radial Laplacian,
\begin{equation}
    \big\langle \grad^2 \field \big\rangle_\mathrm{shell}
    \; \approx \;
    \frac{1}{\rad^2} \partial_\rad\!\big( \rad^2 \partial_\rad \field \big) .
    \label{eq:radial-laplacian-r}
\end{equation}
This deterministic coupling survives even in the strict large-shell limit
in which the noise correlation of the previous subsection vanishes: it is
the leading $\Or(\sigma^2)$ term in the gradient expansion and carries no
extra CLT suppression.

% ======================================================================
\section{Comoving shell geometry}
\label{sec:geometry}
% ======================================================================

Label shells by the step $i$ at which they detach from the active core.
With scale factor normalized so that $\sfactor(\Ninit) = \sfactor_0$, the
outer radius of shell $i$ is the comoving Hubble radius at step $i$,
\begin{equation}
    \rad_i = \frac{1}{\sfactor(i)\Hub(i)} ,
\end{equation}
and the shell width is $\Delta_i = \rad_i - \rad_{i+1}$. Because
$\d(\aH[^{-1}])/\d\Ny = -(1-\epsSR)\aH[^{-1}]$, the shell widths are not
uniform in comoving coordinates: outer shells (small $i$, born early, when
the comoving Hubble radius was larger) are wider than inner shells. The
comoving population of shell $i$ is fixed for all time once the shell is
created; what changes at later steps is only the field value it carries.

\begin{panel}
A discrete implementation along these lines --- a fixed number of
`active' stochastic shells (small enough that the CLT suppression
\eqref{eq:CLT-suppression} has not yet made the noise negligible), coupled
to an outer set of shells evolving under the deterministic gradient-coupled
equation only --- is a viable, and probably cheaper, alternative to the
continuum construction of Sections~\ref{sec:coordinate}--\ref{sec:collocation-scheme}.
The discrete and continuum treatments are two numerical schemes for the
\emph{same} underlying mean-field action, not different physical models. We
record only the continuum route in detail here, since it is the one
adopted for the first implementation, but the discrete route remains
available as a cross-check and may in the end prove more tractable; see
Section~\ref{sec:open-issues} and the companion planning document.
\end{panel}

% ======================================================================
\section{The onion coordinate $\yco$}
\label{sec:coordinate}
% ======================================================================

\subsection{Definition}

We wish to describe the geometry exterior to the horizon
volume that forms the core of the instanton.
First, determine
the current horizon using the core's own local field
values,
\begin{equation}
    \rH(\Ny) \equiv \frac{1}{a(\Ny)\,\Hub\big[\field(1,\Ny),\mom(1,\Ny)\big]} ,
    \label{eq:rH-definition}
\end{equation}
the same local Friedmann-equation evaluation used throughout this
construction, now applied at the domain's own inner edge. The outer edge
of the assembling overdensity is at a fixed comoving radius,
\begin{equation}
    \rout \equiv \frac{1+\alpha}{\aHo} ,
    \qquad \alpha \geq 0 ,
    \label{eq:rout-alpha}
\end{equation}
a small, controlled deformation of the naive choice $\rout=1/\aHo$ used in
an earlier draft, for reasons given in the panel below. We define the
radial variable $s \equiv \ln(\rad/\rout)$ --- motivated by the density
profile computed in the gradient-free approximation being close to a
power law in $\rad$, hence close to linear (and smooth) in $s$ --- and
write
\begin{equation}
    \Deltas(\Ny) \equiv \ln\!\left(\frac{\rout}{\rH(\Ny)}\right) ,
    \qquad
    \yco \equiv -\frac{2s}{\Deltas(\Ny)} - 1 ,
    \qquad
    \rad(\yco,\Ny) = \rout \exp\!\left[-\frac{(\yco+1)\Deltas(\Ny)}{2}\right] .
    \label{eq:y-definition}
\end{equation}
so that $\yco=-1$ is the fixed outer edge $\rout$ (which
will turn out to satisfy
a Dirichlet boundary condition),
and $\yco=+1$ is the current horizon $\rH(\Ny)$ (which
will turn out to satisfy
a Neumann boundary condition). The
domain is therefore $\yco\in[-1,1]$ throughout, matching the natural range
of Legendre/Chebyshev collocation nodes (Section~\ref{sec:collocation-basis}).
Since $\rH(\Ninit)=1/\aHo$ exactly (the core has not yet deviated from the
noiseless trajectory at $\Ninit$),
\begin{equation}
    \Deltas(\Ninit) = \ln(1+\alpha) ,
    \label{eq:Deltas-init}
\end{equation}
strictly positive for $\alpha>0$: the resolved domain starts with a small
but finite log-radial extent, rather than as a single point.

\begin{panel}
\semibold{Why $\alpha>0$: the naive $\rout=1/\aHo$ choice is singular at $\Ninit$.}
With $\alpha=0$, $\Deltas(\Ninit)=0$ exactly, and the Laplacian coefficients
$4/\Deltas^2$, $2/\Deltas$ (eq. 4.6), the advection coefficient $A\propto1/\Deltas$,
and the noise dilution $\ncount\propto\Deltas$ (Section~\ref{sec:noise-y},
so $\Dnoise{ij}\propto1/\Deltas$) all diverge there. A Frobenius-type check
of the forward sector ($\delta\field\equiv\field-\phinl=\Deltas(\Ny)h(\yco)+\Or(\Deltas^2)$)
shows the leading $1/\Deltas$ divergences in the gradient and advection
terms cancel against each other, consistent with the singularity being an
artefact of rescaling a shrinking physical domain onto the fixed range
$[-1,1]$ (transparent if one works instead in $s$, which has no such
prefactor and simply ranges over the perfectly regular, if shrinking,
interval $[-\Deltas(\Ny),0]$) rather than a genuine breakdown of the
physics. But the response-field sector, which carries boundary data from
$\Nfinal$ and couples through the equally-divergent $\Dnoise{ij}\propto1/\Deltas$
noise terms, has not been checked with the same rigor, and there is no
guarantee its divergences cancel as cleanly.

Rather than complete that analysis, decouple $\rout$ from $\rH(\Ninit)$ by
the small amount $\alpha$ in \eqref{eq:rout-alpha}. Since
$\Deltas_\alpha(\Ny) = \ln(1+\alpha) + \Deltas_0(\Ny)$ (with $\Deltas_0$ the
original, unregularized quantity) is just an additive constant shift, every
divergent coefficient above is regularized identically, with no change to
any functional form --- and since the computed profiles already have
$\Deltas_0(\Nfinal)\sim9$--$21$, the shift is negligible almost immediately
after $\Ninit$ for any small $\alpha$, so late-time behaviour is
essentially unaffected. The physical content added is minimal: the thin
shell between $\rH(\Ninit)$ and the new $\rout$ is asserted to sit exactly
on the noiseless trajectory, extending an approximation already made
everywhere beyond $\rout$ slightly further in, not introducing a new kind
of assumption. Smaller $\alpha$ makes the (now finite) coefficients larger
near $\Ninit$, so there is a real stiffness trade-off in the choice; the
recommended approach is to run at a spread of $\alpha$ values (e.g.\
geometrically spaced) and check for the expected negligible sensitivity,
extrapolating to $\alpha\to0$ if any residual dependence remains, rather
than committing to a single value a priori.
\end{panel}

\begin{panel}
\semibold{This was the previously-rejected alternative --- now viable.} An
earlier version of this document considered exactly this construction
(anchoring the inner edge to the current horizon) and rejected it, on the
grounds that the resulting $\Ny$-dependent operator has no time-independent
eigenbasis, forcing re-diagonalization at every step. That objection is
specific to an \emph{explicit spectral method} (solving for coefficients
in an eigenfunction expansion); it does not apply to a \emph{collocation}
method, where the coupling introduced by an $\Ny$-dependent operator is no
worse in kind than the coupling already present from the potential and
noise terms (Section~\ref{sec:collocation-scheme}). Having settled on
collocation for other reasons, the original objection to this coordinate
no longer holds, and it directly resolves the sub-horizon-noise problem
(previously Section~13.1 in an earlier draft) by construction: there is no
region interior to the current horizon left in the domain to source
spurious independent fluctuations.
\end{panel}

\subsection{The advection term}

This
map is genuinely time-dependent at fixed $\rad$, since $\Deltas(\Ny)$
evolves. Writing $s=-(\yco+1)\Deltas(\Ny)/2$ at fixed $\rad$ (i.e.\ fixed
$s$) and differentiating,
\begin{equation}
    \left.\frac{\partial \yco}{\partial \Ny}\right|_\rad
    =
    -\frac{(\yco+1)}{\Deltas(\Ny)}\,\frac{\d\Deltas(\Ny)}{\d\Ny} ,
    \qquad
    \frac{\d\Deltas(\Ny)}{\d\Ny} = -\frac{\d\ln\rH(\Ny)}{\d\Ny} = 1-\epscore(\Ny) ,
    \label{eq:Deltas-derivative}
\end{equation}
using the standard identity $\d\ln(\aH)/\d\Ny = 1-\epsSR$ applied to
$\rH(\Ny)=1/\aH$, evaluated along the core's own trajectory --- valid for
\emph{any} trajectory, since it is essentially the definition of
$\epsSR$, not a background-specific fact. Define the advection coefficient
\begin{equation}
    \advcoef(\yco,\Ny) \equiv \frac{\yco+1}{\Deltas(\Ny)}\big(1-\epscore(\Ny)\big) ,
\end{equation}
so that
\begin{equation}
    \left.\frac{\partial}{\partial \Ny}\right|_\rad
    =
    \left.\frac{\partial}{\partial \Ny}\right|_\yco
    -
    \advcoef(\yco,\Ny)\,\partial_\yco .
    \label{eq:advection-operator}
\end{equation}
$\advcoef$ vanishes at $\yco=-1$ (the fixed outer edge is correctly
undisturbed) and is largest at $\yco=+1$ (the moving horizon).

\subsection{The gradient operator in $\yco$}
\label{sec:gradient-operator}

Using the standard log-radial identity $\nabla^2\field = \rad^{-2}(\field_{ss}+\field_s)$
and $\partial_s = -2 [ \Deltas(\Ny) ]^{-1} \partial_\yco$ (constant at fixed
$\Ny$, so no further chain-rule correction is needed for the second
derivative),
\begin{equation}
    \nabla^2\field
    =
    \rout^{-2}\,\e{\Deltas(\Ny)}\,\e{\Deltas(\Ny)\,\yco}
    \left[\frac{4}{\Deltas(\Ny)^2}\partial_\yco^2\field - \frac{2}{\Deltas(\Ny)}\partial_\yco\field\right]
    \equiv
    \rout^{-2}\,\Lop\field ,
    \label{eq:Lop-definition}
\end{equation}
using $1/\rad(\yco,\Ny)^2 = \rout^{-2}\,\e{\Deltas(\Ny)}\e{\Deltas(\Ny)\yco}$
from \eqref{eq:y-definition} --- note this uses $\rout$ itself, \emph{not}
$\aHo$: with $\alpha>0$ (eq. 4.1a) these are no longer reciprocals of each
other, $\rout=(1+\alpha)/\aHo$, and using $\aHo^2$ here (as an earlier
draft did, implicitly assuming $\alpha=0$) would be wrong by a factor of
$(1+\alpha)^2$. This operator
carries an explicit first-derivative (drift) term and depends on $\Ny$
through $\Deltas(\Ny)$, not just through an overall $\yco$-independent
prefactor. The gradient term in the $\mom$ equation of motion is
\begin{equation}
    \frac{1}{\aHloc^2}\big\langle\grad^2\field\big\rangle
    =
    \frac{1}{\rout^2\aHloc^2}\,\Lop\field ,
    \label{eq:gradient-term-y}
\end{equation}
with $\aHloc$ the same local quantity as before (Section~\ref{sec:noise-y}).

\subsubsection{Self-adjoint measure}

Writing $\Lop$ in $p(\yco)\field''+q(\yco)\field'$ form (dropping the
$\yco$-independent prefactor $\rout^{-2}\e{\Deltas}\cdot4/\Deltas^2$ for
clarity, and absorbing it back at the end), $p=\e{\Deltas\yco}$,
$q=-\tfrac12\Deltas\,\e{\Deltas\yco}$. The Sturm--Liouville reduction
$\measure^{-1}(\measure p\,\field')'$ requires $\measure'/\measure=(q-p')/p$;
here $(q-p')/p = -\tfrac32\Deltas(\Ny)$, a genuine constant (not a function
of $\yco$), so the reduction is exact and closed-form:
\begin{equation}
    \measure(\yco,\Ny) = \e{-\frac32\Deltas(\Ny)\,\yco} ,
    \qquad
    \Lop\field
    \;\propto\;
    \e{\Deltas(\Ny)\yco}\left(\e{-\frac32\Deltas(\Ny)\yco}\right)^{-1}\frac{\d}{\d\yco}\!\left(\e{-\frac12\Deltas(\Ny)\yco}\,\frac{\d\field}{\d\yco}\right) .
    \label{eq:self-adjoint-measure}
\end{equation}
As a consistency check, $\rad(\yco,\Ny)^3 \propto \e{-\frac32\Deltas(\Ny)\yco}$
--- $\measure(\yco,\Ny)$ is, up to normalization, exactly the physical
comoving volume measure, as it must be: $\nabla^2$ is self-adjoint with
respect to $\rad^2\,\d\rad$ in \emph{any} radial parametrization, by
Green's theorem, so recovering it here from the Sturm--Liouville reduction
is a check on the algebra rather than a new input.

\begin{panel}
\semibold{Flat-measure (Schrödinger normal) form: explored and set aside.}
The further Liouville substitution that removes $q(\yco)$ \emph{and}
brings the measure to flat $\d\yco$ (rather than the weighted
$\measure(\yco,\Ny)\,\d\yco$ above) was worked through explicitly. It
succeeds algebraically --- the residual potential term vanishes exactly,
an exactly-solvable special case, since the exponential coefficients here
play the same role as in an Euler-type ODE --- but the new independent
variable $x(\yco)=\tfrac{2}{\Deltas(\Ny)}(1-\e{-\Deltas(\Ny)\yco/2})$ maps
$\yco\in[-1,1]$ onto a badly asymmetric interval whose width grows like
$\e{\Deltas(\Ny)/2}$: for the wide ($\dNstar=13$) profile already computed
in the gradient-free approximation, $\Deltas\approx20.7$ gives a domain
width $\sim3\times10^3$; even the narrower ($\dNstar\approx1.5$) case gives
$\Deltas\approx9.2$ and a width $\sim22$. This reintroduces, in the
transformed coordinate, essentially the same exponential compression the
logarithmic radial variable was adopted to remove. \textbf{Conclusion:}
retain the \emph{weighted} measure \eqref{eq:self-adjoint-measure} and
build the self-adjoint discrete structure (Section~\ref{sec:collocation-scheme})
directly for it, rather than transforming to a flat measure.
\end{panel}

This removes the previous $\S4.4$ exponential-mode-resolution concern
structurally, not by mitigation: the logarithmic radial variable is
precisely what turns the exponentially-compressed near-horizon structure
(in linear $\rad$) into $\Or(1)$ structure in $\yco$ --- consistent with
the density profile already computed in the gradient-free approximation
being close to a power law in $\rad$, hence smooth and close to linear in
$s=\ln(\rad/\rout)$. A first implementation should confirm this
empirically once collocation-point convergence can be checked, but there
is no longer a structural reason to expect the mode/point count needed to
grow exponentially with elapsed $e$-folds.

% ======================================================================
\section{Noise statistics in $\yco$ coordinates}
\label{sec:noise-y}
% ======================================================================


The number of coarse-graining (current-Hubble-volume) patches per unit
$\yco$-shell can be computed purely in comoving terms. From
\eqref{eq:y-definition}, $\d\rad/\d\yco = -\tfrac12\Deltas(\Ny)\,\rad(\yco,\Ny)$
at fixed $\Ny$, so a shell of coordinate width $\d\yco$ has comoving volume
$4\pi\rad^2|\d\rad| = 2\pi\Deltas(\Ny)\,\rad(\yco,\Ny)^3\,\d\yco$. The
comoving volume of the current Hubble patch \emph{local to that shell} ---
which is what matters, since noise is a spatially local process --- is
$\tfrac43\pi/\aHloc^3$, using the same local $\aHloc$ as the gradient term
(Section~\ref{sec:gradient-operator}). The number of patches per unit
$\yco$ is the ratio of these volumes,
\begin{equation}
    \ncount(\yco,\Ny)
    =
    \frac32\Deltas(\Ny)\,\big[\rad(\yco,\Ny)\,\aHloc\big]^3
    =
    \frac32\Deltas(\Ny)\,\big[\rout\,\aHloc\big]^3 \, \e{-\frac32(\yco+1)\Deltas(\Ny)} ,
    \label{eq:ncount}
\end{equation}
using $\rad(\yco,\Ny) = \rout\,\e{-(\yco+1)\Deltas(\Ny)/2}$ directly from
\eqref{eq:y-definition} --- note this is written in terms of $\rout$, not
$\aHo$: since $\rout=(1+\alpha)/\aHo$, these differ once $\alpha>0$, and
the combination $\rout\aHloc$ is the one that appears without extra
factors of $(1+\alpha)$. As before, the $\aH$ inside $\aHloc$ is the local
quantity used throughout Section~\ref{sec:instanton-eqs}, while $\rout$
retains its fixed, coordinate-defining role. At the outer edge
($\yco=-1$, where $\field=\phinl$ is pinned), $\aHloc\to a(\Ny)\Hub[\phinl(\Ny),\pinl(\Ny)]$,
generically $\gg\aHo$ (hence $\rout\aHloc\gg1$, since $\rout=\Or(1/\aHo)$)
at late $\Ny$, so $\ncount$ is large there (many patches, as expected). At
the horizon itself ($\yco=+1$), $\aHloc$ is by \eqref{eq:rH-definition}
exactly $1/\rH(\Ny)$, so $\rout\aHloc = \rout/\rH(\Ny) = \e{\Deltas(\Ny)}$
directly from the definition of $\Deltas(\Ny)$, and the two exponential
factors cancel: $\ncount(1,\Ny) = \tfrac32\Deltas(\Ny)$
--- an $\Or(10)$, not $\Or(1)$, quantity for the profiles already computed
($\Deltas\sim9$--$21$ depending on $\dNstar$): the exterior-only domain no
longer approaches the single-patch limit $\ncount\to1$ at its inner edge,
because the excised sub-horizon interior --- not the last resolved shell
--- is where that single remaining patch now lives, outside the domain
entirely.

\begin{panel}
\semibold{$\ncount(\yco,\Ny)$ as a density, not a shell count.}
\eqref{eq:ncount} is a \emph{density} $\d N_{\mathrm{patches}}/\d\yco$, not
the (dimensionless) patch count in a shell of finite width: the count in a
shell of coordinate width $\d\yco$ is $\ncount(\yco,\Ny)\,\d\yco$, and it
is the density that enters $\Dnoise{\field}=\Pspec{\field\field}/\ncount$
below, consistently with the $\delta(\yco-\yco')$ noise correlation of
\eqref{eq:langevin-y}, which requires a density on the right-hand side
rather than a quantity that vanishes as $\d\yco\to0$.
\end{panel}

\begin{panel}
\semibold{Geometric dilution versus local amplitude.} $\ncount(\yco,\Ny)$
above is \emph{not} purely a coordinate/gauge quantity, unlike $\aHo,\rH(\Ny)$
as used to define $\yco$ itself: it uses $\aHloc$, the same local $aH$ as
the gradient term, because ``how big is the current Hubble patch at this
shell'' is a genuinely local physical statement, not a bookkeeping choice
about the comoving coordinate map. The single-patch \emph{amplitude}
entering the average, $\Dcm{ij}(\field,\mom)$, is a separate object with
its own local dependence: the codebase already treats the diffusion matrix
as a local function of the phase-space point --
\texttt{DiffusionModel.D\_matrix(phi,pi)} returns $(\Dcm{11},\Dcm{12},\Dcm{22})$
evaluated at $(\field,\mom)$, exactly as \texttt{AbstractPotential} does for
$V$, $H^2$, $\epsilon_1$ -- and the existing single-trajectory equations of
motion (\texttt{ComputeTargets/FullInstanton.py}) use it that way,
e.g.\ $\dot\field = \mom + 2\Dcm{11}\rfield + 2\Dcm{12}\rmom$.
\end{panel}

Using $\Pspec{fg}\atloc = 2\,\Dcm{fg}\atloc$ (the identification implicit in
$\Pspec{}/2 = H^2/8\pi^2$ in the slow-roll, decoupled-diffusion case), the
correctly local shell-averaged noise variance/covariance densities are
\begin{equation}
    \Dnoise{\field}(\yco,\Ny) \equiv \frac{2\,\Dcm{11}\big(\field(\yco,\Ny),\mom(\yco,\Ny)\big)}{\ncount(\yco,\Ny)} ,
    \qquad
    \Dnoise{\mom}(\yco,\Ny) \equiv \frac{2\,\Dcm{22}\big(\field(\yco,\Ny),\mom(\yco,\Ny)\big)}{\ncount(\yco,\Ny)} ,
    \label{eq:Dnoise-diag}
\end{equation}
together with the cross term,
\begin{equation}
    \Dnoise{\field\mom}(\yco,\Ny) \equiv \frac{2\,\Dcm{12}\big(\field(\yco,\Ny),\mom(\yco,\Ny)\big)}{\ncount(\yco,\Ny)} .
    \label{eq:Dnoise-cross}
\end{equation}
All three are evaluated pointwise on the same $\yco$-grid, using the same
reconstructed $\field(\yco,\Ny)$, $\mom(\yco,\Ny)$, by the same procedure
used for every other nonlinear, solution-dependent coefficient in this
scheme (Section~\ref{sec:collocation-scheme}). Since $\ncount$ no longer
approaches $1$ anywhere in the resolved domain, treating
$\Dnoise{\field}(\yco,\Ny)$ as a pointwise, $\delta(\yco-\yco')$-correlated
density is on safer footing throughout $\yco\in[-1,1]$ than it was for the
previous, interior-covering coordinate.
% ======================================================================
\section{The Langevin system in $(\yco,\Ny)$}
\label{sec:langevin-y}
% ======================================================================

Collecting the results of Sections~\ref{sec:coordinate}--\ref{sec:noise-y},
the mean-field Langevin system is
\begin{subequations}
\label{eq:langevin-y}
\begin{align}
    \frac{\partial \field}{\partial \Ny}
    & =
    \mom + \advcoef(\yco,\Ny)\,\partial_\yco\field + \noise{\field}(\yco,\Ny) ,
    \\
    \frac{\partial \mom}{\partial \Ny}
    & =
    -\big(3-\epsSR\atloc\big)\mom - \frac{\Vpot'(\field)}{\Hub^2\atloc}
    + \frac{1}{\rout^2\aHloc^2} \Lop \field
    + \advcoef(\yco,\Ny)\,\partial_\yco\mom
    + \noise{\mom}(\yco,\Ny) ,
\end{align}
\end{subequations}
where the derivatives on the left are now at fixed $\yco$
(\eqref{eq:advection-operator} converts the original fixed-$\rad$ equations),
and the advection term $\advcoef\,\partial_\yco(\cdot)$ --- absent for the
previous coordinate, present here because the map itself is
$\Ny$-dependent (Section~\ref{sec:coordinate}) --- appears in \emph{both}
equations, since both were originally written at fixed $\rad$. As before,
$\Hub^2\atloc$ and $\epsSR\atloc$ denote $H_{\mathrm{sq}}[\field(\yco,\Ny),\mom(\yco,\Ny)]$
and $\epsilon[\field(\yco,\Ny),\mom(\yco,\Ny)]$, the same local
Friedmann-equation functions used elsewhere, distinct from the genuinely
background $\aHo,\rH(\Ny)$ used to define $\yco$ itself and
$\ncount(\yco,\Ny)$. On the fixed rectangular domain $\yco\in[-1,1]$,
$\Ny\in[\Ninit,\Nfinal]$, the noises have covariance
\begin{equation}
    \big\langle \noise{\field}(\yco,\Ny)\noise{\field}(\yco',\Ny')\big\rangle
    =
    \Dnoise{\field}(\yco,\Ny)\,\delta_D(\Ny-\Ny')\,\delta(\yco-\yco') ,
\end{equation}
similarly for $\big\langle\noise{\mom}\noise{\mom}\big\rangle$ with
$\Dnoise{\mom}$, and
\begin{equation}
    \big\langle \noise{\field}(\yco,\Ny)\noise{\mom}(\yco',\Ny')\big\rangle
    =
    \Dnoise{\field\mom}(\yco,\Ny)\,\delta_D(\Ny-\Ny')\,\delta(\yco-\yco') ,
\end{equation}
where the $\delta(\yco-\yco')$ reflects the result of
Section~\ref{sec:shell-averaging} that inter-shell noise correlations are
subleading once the CLT suppression has been accounted for. Unlike the
previous coordinate, there is no longer any sub-horizon region within the
domain to worry about: $\yco=+1$ \emph{is} the current horizon by
construction, so the domain boundary and the physical boundary between
``many independent patches'' and ``at most one patch's worth of freedom''
coincide exactly, at every $\Ny$.

% ======================================================================
\section{The MSR action}
\label{sec:msr-action}
% ======================================================================

\subsection{Choice of measure}

A generic MSR action for the Langevin system \eqref{eq:langevin-y} could be
written with a flat $\d\yco$ measure. However, the physically correct
measure is the one that makes $\Lop$ self-adjoint. This is
\begin{equation}
    \measure(\yco,\Ny) = \e{-\frac32\Deltas(\Ny)\,\yco} ,
\end{equation}
which, per the discussion following \eqref{eq:self-adjoint-measure}, is
(up to normalization) exactly the physical comoving volume element in this
coordinate.
We write the MSR action with this measure,
introducing response fields $\rfield$, $\rmom$ conjugate to $\field$,
$\mom$, and including the advection terms from \eqref{eq:advection-operator}
in both bracketed equations of motion:
\begin{multline}
    S
    =
    \int_{\Ninit}^{\Nfinal}\! \d\Ny \int_{-1}^1\! \measure(\yco,\Ny)\,\d\yco \; \Bigg[
    \rfield \Big( \dot\field - \mom - \advcoef\,\partial_\yco\field \Big)
    \\
    +
    \rmom \left( \dot\mom + \big(3-\epsSR\atloc\big)\mom + \frac{\Vpot'(\field)}{\Hub^2\atloc}
    - \frac{1}{\rout^2\aHloc^2}\Lop\field - \advcoef\,\partial_\yco\mom \right)
    -
    \frac{\Dnoise{\field}}{2}\rfield^2
    -
    \Dnoise{\field\mom}\,\rfield\,\rmom
    -
    \frac{\Dnoise{\mom}}{2}\rmom^2
    \Bigg] ,
    \label{eq:msr-action}
\end{multline}
where an overdot denotes $\partial/\partial\Ny$ at fixed $\yco$, and the
cross term $-\Dnoise{\field\mom}\rfield\rmom$ has coefficient $1$, not
$\tfrac12$, being the off-diagonal part of $-\tfrac12\rfield_i\Dnoise{ij}\rfield_j$.

\begin{panel}
\semibold{Check of self-adjointness, with the drift term.} Writing $\Lop$
(dropping the $\Ny$-dependent overall constant $\rout^{-2}\e{\Deltas}\cdot4/\Deltas^2$
from \eqref{eq:Lop-definition}, which factors out of the check) as
$p(\yco)\,Q''+q(\yco)\,Q'$ with $p(\yco)=e^{\Deltas\yco}$,
$q(\yco)=-\tfrac12\Deltas e^{\Deltas\yco}$, and writing
$\weightfn(\yco)\equiv \measure(\yco,\Ny)\,p(\yco) = e^{-\Deltas\yco/2}$
(the product of the self-adjoint measure with the coefficient of $Q''$,
\emph{not} its inverse: $\measure=\weightfn^3$, and $\weightfn'=-\tfrac12\Deltas\,\weightfn$),
\begin{align*}
    \int \measure\,\d\yco\; P\,\Lop Q
    & \propto
    \int \weightfn\,\d\yco\;P\,Q'' \;-\; \frac{\Deltas}{2}\int\weightfn\,\d\yco\;P\,Q'
    \\
    & =
    \Big[\weightfn P Q'\Big]_{\mathrm{bdry}} - \int \weightfn'\,P\,Q'\,\d\yco - \int\weightfn\,P'\,Q'\,\d\yco - \frac{\Deltas}{2}\int\weightfn\,P\,Q'\,\d\yco
    \\
    & =
    \Big[\weightfn P Q'\Big]_{\mathrm{bdry}} - \int \weightfn\, P'\,Q'\,\d\yco ,
\end{align*}
using $\weightfn'=-\tfrac12\Deltas\weightfn$ to cancel the two
$\weightfn'PQ'$-type terms exactly. The bulk term $-\int\weightfn P'Q'\,\d\yco$
is manifestly symmetric in $P,Q$; self-adjointness therefore requires only
that the boundary term $\big[\weightfn(PQ'-QP')\big]_{\mathrm{bdry}}$
vanish.
At $\yco=+1$, the regularity conditions $\partial_\yco\field=0$,
$\partial_\yco\rmom=0$ (Section~\ref{sec:bcs}) give $P'=Q'=0$ there,
killing the boundary term directly.
At $\yco=-1$, $\rfield=\rmom=0$ (no response where the outer data
is pinned), giving $P=Q=0$ there.
\end{panel}

% ======================================================================
\section{Instanton equations}
\label{sec:instanton-eqs}
% ======================================================================

Extremizing \eqref{eq:msr-action} with respect to $\rfield,\rmom,\field,\mom$
gives the coupled instanton system
\begin{subequations}
\label{eq:instanton-eqs}
\begin{align}
    \dot\field
    & =
    \mom + \advcoef\,\partial_\yco\field + \Dnoise{\field}\,\rfield + \Dnoise{\field\mom}\,\rmom ,
    \label{eq:inst-phi}
    \\
    \dot\mom
    & =
    -\big(3-\epsSR\atloc\big)\mom - \frac{\Vpot'(\field)}{\Hub^2\atloc}
    + \frac{1}{\rout^2\aHloc^2}\Lop\field
    + \advcoef\,\partial_\yco\mom
    + \Dnoise{\mom}\,\rmom + \Dnoise{\field\mom}\,\rfield ,
    \label{eq:inst-pi}
    \\
    \dot{\rfield}
    & =
    -\frac{\Vpot''(\field)}{\Hub^2\atloc}\rmom
    + \frac{1}{\rout^2\aHloc^2}\Lop\rmom
    - \frac{1}{\measure}\partial_\yco\big(\measure\,\advcoef\,\rfield\big) ,
    \label{eq:inst-rphi}
    \\
    \dot{\rmom}
    & =
    \rfield - \big(3-\epsSR\atloc\big)\rmom
    - \frac{1}{\measure}\partial_\yco\big(\measure\,\advcoef\,\rmom\big) .
    \label{eq:inst-rpi}
\end{align}
\end{subequations}
The response fields $\rfield,\rmom$ take the place of the noises
$\noise{\field}, \noise{\mom}$ in \eqref{eq:inst-phi}--\eqref{eq:inst-pi},
with the cross term $\Dnoise{\field\mom}$ coupling each
response field to the \emph{other} noise channel. Both
$\dot\field$ and $\dot\mom$ pick up the
advection term $\advcoef\partial_\yco(\cdot)$ directly (from
\eqref{eq:advection-operator}), and \emph{both} response-field equations
pick up a term $-\measure^{-1}\partial_\yco(\measure\,\advcoef\,X)$
built from that \emph{same} response field $X\in\{\rfield,\rmom\}$ ---
the standard MSR result that forward advection $+\advcoef\partial_\yco$
becomes its negative adjoint (w.r.t.\ $\measure$) in the backward/response
equations, rather than mixing $\rfield$ and $\rmom$ the way the potential
and gradient terms do. Expanding,
\begin{equation}
    \frac{1}{\measure}\partial_\yco\big(\measure\,\advcoef\,X\big)
    =
    \advcoef\,\partial_\yco X + X\,(1-\epscore)\left[\frac{1}{\Deltas} - \frac32(\yco+1)\right] ,
\end{equation}
using $\partial_\yco\advcoef = (1-\epscore)/\Deltas$ and
$\partial_\yco\measure/\measure = -\tfrac32\Deltas$, both constants at
fixed $\Ny$.

\begin{panel}
\semibold{Status of this derivation.} The advection/response-field coupling
above was worked out in this pass rather than carried over from earlier,
independently-checked material, and has not had the same scrutiny as the
rest of this section. The structure matches the standard MSR result for
advection-diffusion (forward $+A\partial_\yco$, backward
$-\measure^{-1}\partial_\yco(\measure A\,\cdot)$), which is a useful
consistency check but not a substitute for redoing the derivation
independently before relying on it numerically.
\end{panel}

\begin{panel}
\semibold{What is \emph{not} corrected here.} Equations \eqref{eq:inst-rphi}--\eqref{eq:inst-rpi}
are obtained by varying \eqref{eq:msr-action} with respect to $\field,\mom$.
Since $\Hub^2\atloc$, $\epsSR\atloc$, $\Dnoise{ij}\atloc$, and now
$\advcoef(\yco,\Ny)$ (through $\epscore(\Ny)$ and $\Deltas(\Ny)$, both
built from the core values $\field(1,\Ny),\mom(1,\Ny)$) are local functions
of $\field,\mom$, a fully consistent variation would generate additional
terms proportional to $\partial\Hub^2/\partial\field$,
$\partial\Dnoise{ij}/\partial\field$, $\partial\advcoef/\partial\field$,
etc., alongside the $\Vpot''(\field)$ term already present. These are
dropped here, consistently with the fact that \texttt{FullInstanton}'s
\texttt{bwd\_rhs} does not include the analogous $H^2,\epsSR$ feedback
terms either (only the $V''$ term appears): a higher-order effect in the
same sense as other terms already truncated in the gradient expansion, not
a new omission. The gradient term needs the same truncation applied
slightly more carefully, since $\aHloc$ sits \emph{inside} $\Lop$'s
argument: treat $\aHloc$ (and now $\advcoef$) as frozen --- evaluated at
the current reconstructed $\field,\mom$, but not differentiated through
--- when varying to obtain \eqref{eq:inst-rphi}--\eqref{eq:inst-rpi}. With
that understood, the self-adjointness check above guarantees the stated
structure only \emph{at this truncation}, not as an exact consequence of
$\Lop$'s self-adjointness once its coefficients are $\yco$-dependent.
\end{panel}

\section{Boundary conditions}
\label{sec:bcs}
% ======================================================================

\subsection{Spatial boundary conditions}

\para{At $\yco=-1$ (fixed outer edge)} The field is pinned to the known
noiseless trajectory,
\begin{equation}
    \field(-1,\Ny) = \phinl(\Ny) ,
    \qquad
    \mom(-1,\Ny) = \pinl(\Ny) ,
    \label{eq:bc-outer}
\end{equation}
and correspondingly $\rfield(-1,\Ny) = \rmom(-1,\Ny) = 0$ (no response
where the field is pinned).

\para{At $\yco=+1$ (current horizon)} We impose
\begin{equation}
    \partial_\yco \field \big|_{\yco=1} = 0 ,
    \qquad
    \partial_\yco \rmom \big|_{\yco=1} = 0 .
    \label{eq:bc-core}
\end{equation}
The reason for imposing this condition
is the physical
content of excising the sub-horizon interior in the first place --- that
it carries no independent structure of its own, so the field attaches
smoothly to whatever lies beyond the domain, with no preferred direction
singled out at the point of attachment. This is a physical assumption
restated as a smoothness condition, rather than a geometric necessity;
compare Section~\ref{sec:coordinate}. As before, $\field(1,\Ny)$ itself
--- read off as the trajectory --- is left free by this condition and is
an output of the solution, not an input.

\begin{panel}
\semibold{A rejected alternative.} One might instead impose \emph{both}
Dirichlet and Neumann data at $\yco=-1$,
$\field(-1,\Ny)=\phinl(\Ny)$ and $\partial_\yco\field(-1,\Ny)=0$, leaving
no condition at $\yco=+1$, so that the presence or absence of a cusp at
the horizon becomes an output rather than an input. This is a Cauchy
problem in $\yco$ for what is, at fixed $\Ny$, an elliptic operator
($\Lop$), and Cauchy problems for elliptic operators are generically
ill-posed in the sense of Hadamard: small perturbations of the data at
$\yco=-1$ can grow uncontrollably on integration towards $\yco=+1$. Mixed
conditions at the two distinct ends --- Dirichlet at $\yco=-1$, Neumann at
$\yco=+1$ --- is the well-posed configuration, exactly as before.
\end{panel}

\subsection{Temporal boundary conditions}

\begin{align}
    \field(\yco,\Ninit) & = \phiinit ,
    &
    \mom(\yco,\Ninit) & = \mom_{\mathrm{SR,init}} ,
    &
    \text{(all } \yco \text{)}
    \label{eq:bc-init}
    \\
    \field(1,\Nfinal) & = \phiend ,
    & & &
    \label{eq:bc-final-core}
    \\
    \rfield(\yco,\Nfinal) & = 0 ,
    &
    \rmom(\yco,\Nfinal) & = 0 ,
    &
    \text{(}\yco<1\text{, free)}
    \label{eq:bc-final-free}
\end{align}
with the constraint \eqref{eq:bc-final-core} enforced via a Lagrange
multiplier $\lag$ that sources $\rfield(\yco,\Nfinal)$ as a point
contribution at $\yco=1$; see Section~\ref{sec:collocation-scheme} for
the resulting terminal condition, considerably simpler here than in the
mode-expansion scheme it replaces (Section~\ref{sec:collocation-scheme}).
No boundary condition is imposed at $\yco<1$ forcing $\field$ to reach
$\phiend$: we do \emph{not} impose an explicit `freezing' of
$\field(\yco,\Ny)$ once it crosses $\phiend$ (see
Section~\ref{sec:zeta-extraction}).

% ======================================================================
\section{No explicit freezing; extraction of $\zprofile(\yco)$}
\label{sec:zeta-extraction}
% ======================================================================

In the single-trajectory instanton, `end of inflation' is a single event.
Here, different $\yco$ reach $\phiend$ at different times
$\Ny_{\mathrm{stop}}(\yco)$, with $\Ny_{\mathrm{stop}}(-1) = $ (the noiseless
crossing time) and $\Ny_{\mathrm{stop}}(1) = \Nfinal$. Treating
$\Ny_{\mathrm{stop}}(\yco)$ as an a priori stopping boundary makes this a
free-boundary problem: the domain of \eqref{eq:instanton-eqs} would have a
curved upper edge that is itself part of the solution.

We avoid this by \emph{not} imposing any stopping or freezing condition:
the system \eqref{eq:instanton-eqs} is integrated over the full rectangle
$\yco\in[-1,1]$, $\Ny\in[\Ninit,\Nfinal]$, with the single core condition
\eqref{eq:bc-final-core}, and $\zprofile(\yco)$ is read off as a
\emph{post-processing} step applied to the resulting $\field(\yco,\Nfinal),
\mom(\yco,\Nfinal)$.

\par\vspace{0.5em}\noindent
\semibold{Scale assignment: density matching, not a crossing scan.}
The instanton can, in principle, leave each shell with its own
isocurvature perturbation relative to the others (small and rapidly
decaying in a single-field model, but not exactly zero), so the physically
correct radius assignment compares surfaces of \emph{fixed energy density}
between the shell and the background, not surfaces of fixed field value on
a fixed-$\Ny$ (flat-gauge-like) slicing -- these differ once $\zprofile$ is
$\Or(1)$, in the same way uniform-density and flat time-slicings differ in
the standard $\delta N$ formalism. Concretely, for each $\yco$:
\begin{enumerate}
    \item Take the reconstructed state $\big(\field(\yco,\Nfinal),\mom(\yco,\Nfinal)\big)$
    at the \emph{shared} final time $\Nfinal$ -- the same $\Ny$ for every
    $\yco$, read directly off the global solution, with no per-shell
    crossing search.
    \item Continue \emph{that shell's own} trajectory noiselessly forward
    from this state until $\epsSR=1$ (genuine end of inflation), giving an
    endpoint $e$-fold $\Ny_{\mathrm{end}}(\yco)$ and an energy density
    $\rho_{\mathrm{end}}(\yco) \equiv \rho\big[\field_{\downarrow}(\yco,\Ny_{\mathrm{end}}),
    \mom_{\downarrow}(\yco,\Ny_{\mathrm{end}})\big]$ there, where the
    $\downarrow$ subscript denotes this per-shell noiseless downflow.
    \item Density-match against the background: find $\Ny_{\mathrm{nl}}(\rho_{\mathrm{end}}(\yco))$,
    the $e$-fold at which the noiseless trajectory itself has energy
    density $\rho_{\mathrm{end}}(\yco)$.
    \item
    \begin{equation}
        \zprofile(\yco) = \Ny_{\mathrm{end}}(\yco) - \Ny_{\mathrm{nl}}\big(\rho_{\mathrm{end}}(\yco)\big) .
        \label{eq:zeta-extraction}
    \end{equation}
\end{enumerate}
This is the same construction \texttt{CompactionFunction} already uses for
the single-trajectory case (downflow to $\epsSR=1$, then match against the
background).
Applying it pointwise in $\yco$ here is the natural generalization,
sidesteps needing to surface-find $\Ny_{\mathrm{stop}}(\yco)$ at all
(addressing the free-boundary concern above by a different route than
originally intended: the downflow does the work that the crossing scan
was trying to do, without needing $\field(\yco,\Nfinal)$ to already be at
or past $\phiend$), and is expected to give monotonically ordered radii
across shells provided $\zprofile$ is not increasing so rapidly with
$\yco$ that it produces a Type~II perturbation.


% ======================================================================
\section{Physical and comoving scale assignment}
\label{sec:scale-assignment}
% ======================================================================

The construction above gives $\zprofile(\yco)$, a genuine profile rather
than the single number the existing single-trajectory pipeline produces.
Turning this into radius and mass estimates for
\texttt{CompactionFunction} (or the equivalent
container for \texttt{GradientCoupledInstanton})
requires distinguishing three different
notions of ``scale'' that the single-trajectory case did not need to
separate, because there a single comoving radius played all three roles
at once.

\subsection{Comoving radius}

$\rad(\yco,\Ny)$ is already known for every $\yco$, directly from the
coordinate map \eqref{eq:y-definition} evaluated at $\Ny=\Nfinal$ once
$\Deltas(\Nfinal)$ has been solved for --- no separate calculation. Being
comoving, it carries none of the metric perturbation: it is a coordinate
label, unperturbed by construction, in the same sense any comoving
coordinate is.

\subsection{Areal radius}

The GR-meaningful radius, and the one entering the standard compaction
function formula (Yoo 2022),
\begin{equation}
    C(\rad) = \frac23\Big[1-\big(1+\rad\,\zprofile'(\rad)\big)^2\Big] ,
    \label{eq:compaction-yoo}
\end{equation}
is the areal radius $R(\rad) = a\,\rad\,e^{\zprofile(\rad)}$: this formula
is derived directly from $\d R/\d\rad = a\,e^{\zprofile}(1+\rad\zprofile')$,
so the areal-radius relation is already built into it. Note
\eqref{eq:compaction-yoo} takes $\zprofile(\rad)$ and
$\zprofile'(\rad)\equiv\d\zprofile/\d\rad$ --- comoving derivatives --- as
input, \emph{not} $\zprofile(R)$: using the areal coordinate here would be
a genuine mismatch, not an equivalent relabelling. Since
$\zprofile(\yco_j)$ is already known at every collocation node and
$\rad(\yco_j,\Nfinal)$ from the coordinate map, $\zprofile'(\rad)$ follows
by the chain rule, $\zprofile'(\rad) = (\d\zprofile/\d\yco)/(\d\rad/\d\yco)$,
using the collocation differentiation matrix for the numerator and the
analytic $\d\rad/\d\yco=-\tfrac12\Deltas(\Nfinal)\rad(\yco,\Nfinal)$ for the
denominator --- no new derivative machinery.
Alternatively, since we will extract $\zeta(r)$ anyway, the
deriative can be computed directly.

$R(\rad)$ is not itself an observable; its role is as the correctly
GR-translated input against which a numerical-relativity-calibrated
threshold $C_{\rm th}$ can be compared, since NR collapse simulations
(Musco, Escriv\'a, Young, and collaborators) are set up and calibrated
directly in areal-radius/compaction-function language. As $\rad$ shrinks
towards the interior, if $\zprofile(\rad)$ grows steeply enough $R(\rad)$
can stop being monotonic in $\rad$ --- the Type~II perturbation regime
already flagged (Section~\ref{sec:zeta-extraction}), corresponding to
$C_{\max}\to2/3$, and the reason the existing pipeline computes $r_{\max}$
and $r_{\rm peak}$ from $C(\rad)$ alone rather than from any
$\bar C$-based estimate sensitive to the $e^{3\zprofile}$ volume
weighting.

\subsection{Physical (present-day) scale}

A third, independent notion: the wavelength a given comoving mode
corresponds to today, needed for the mass/scale axis of the final PBH
mass function. This uses the existing Leach--Liddle scale-matching
equation (astro-ph/0305263), solved \emph{once}, anchored at the fixed
outer edge $\rout$ --- exactly as for the single-trajectory case, since
$\yco=-1$ sits exactly on the noiseless background by construction
(eq. 9.1), giving a physical scale $\rad_{\mathrm{phys,out}}$. Since comoving
wavenumber satisfies $k\propto1/\rad$, every other shell's physical scale
follows by a fixed ratio,
\begin{equation}
    \rphys(\yco) = \frac{\rad(\yco,\Nfinal)}{\rout}\;\rad_{\mathrm{phys,out}} ,
    \label{eq:rphys-ratio}
\end{equation}
with no further per-shell Leach--Liddle solve and no per-shell inversion
of the background expansion history: comoving-ness already fixes the
today's-scale conversion to a single universal constant, common to every
mode regardless of when it crossed the horizon, and that constant is
exactly what solving Leach--Liddle once determines. The genuine
per-shell/isocurvature-sensitive calculation --- the downflow from each
shell's own instanton endpoint to $\epsSR=1$, giving that shell's own
$V_{\rm end,\,downflow}$ and hence its own duration in the sense already
used for $\zprofile(\yco)$ (Section~\ref{sec:zeta-extraction}) --- is
unaffected and answers a different physical question (how long \emph{that
particular shell} takes to finish inflating) from "which comoving scale is
this" (a purely geometric fact, answered once and for all by
\eqref{eq:rphys-ratio}).

\begin{panel}
\textbf{A candidate refinement, considered and rejected.} One might
instead try to assign each shell its own horizon-crossing $e$-fold by
inverting $a(\Ny')\Hub(\Ny')=k[\rad(\yco,\Nfinal)]$ against the
\emph{background} trajectory. This was considered and is unnecessary: it
answers a question (``which $e$-fold did this comoving scale cross the
horizon'') that plays no role once $\rad$ is already comoving, since the
comoving-to-today conversion is a single constant for every mode by
definition, not something that depends on horizon-crossing time
individually. The single global ratio \eqref{eq:rphys-ratio} already
contains everything the inversion would have supplied.
\end{panel}

\subsection{Equivalence check against the discrete (peeling) scheme}

Applying \eqref{eq:rphys-ratio} to the core, $\yco=1$: since
$\rad(1,\Ny)\equiv\rH(\Ny)$ by the definition of the coordinate's inner
edge \eqref{eq:rH-definition}, this gives
$\rad_{\mathrm{phys,core}} = \rad_{\mathrm{phys},H}(\Nfinal)$ directly. In the discrete
peeling construction (Section~\ref{sec:geometry}), the innermost shell at
step $\Nfinal$ has comoving radius $\rad_i=1/[a(\Nfinal)\Hub(\Nfinal)]$ by
the \emph{same} defining formula used for every layer there --- the same
geometric object, playing the same role, so the two schemes assign the
core the same physical scale by construction, not merely in numerical
agreement. This holds exactly as $\alpha\to0$ (Section~\ref{sec:coordinate}),
with $\Or(\alpha)$ corrections at finite $\alpha$, consistent with $\alpha$
being a small, controlled deformation whose effects are expected to
extrapolate away.

\begin{todo}
This equivalence assumes the discrete construction's $\rad_i=1/[a(i)\Hub(i)]$
uses the instanton's own local trajectory (the actual, possibly
noise-driven core), matching how $\rH(\Ny)$ is defined here via
$\Hub[\field(1,\Ny),\mom(1,\Ny)]$. This has not been checked against the
existing implementation of the discrete/peeling scheme; if that code in
fact uses the background trajectory's $a,\Hub$ for $\rad_i$, the two
schemes would agree only in the noiseless limit, not in general, and this
section's equivalence claim would need revisiting.
\end{todo}
% ======================================================================
\section{Collocation basis: Legendre--Gauss--Lobatto nodes}
\label{sec:collocation-basis}
% ======================================================================

Rather than expand in eigenfunctions of $\Lop$ (not available in closed
form once $\Lop$ carries the drift term of Section~\ref{sec:gradient-operator}),
we
represent $\field,\mom,\rfield,\rmom$ by their values on a fixed grid of
collocation points and work directly with differentiation matrices. This
was chosen over Chebyshev collocation specifically so that the
self-adjoint structure of $\Lop$ (Section~\ref{sec:msr-action}) can be
preserved numerically via a summation-by-parts (SBP) operator with a
\emph{diagonal} norm matrix --- the standard construction for
Legendre--Gauss--Lobatto (LGL) nodes, not available for Chebyshev without
retrofitting a non-diagonal norm.

\subsection{Nodes, weights, and differentiation matrices}

The $\modeidx_{\max}+1$ LGL nodes $\{\yco_j\}$ on $[-1,1]$ are $\pm1$
together with the interior roots of $P_{\modeidx_{\max}}'(x)=0$ (the
derivative of the Legendre polynomial of degree $\modeidx_{\max}$),
computed via the eigenvalues of a symmetric tridiagonal Jacobi matrix ---
standard, tabulated, and readily available (e.g.\ the Golub--Welsch-type
routines in \texttt{FastGaussQuadrature.jl}, or the reference Python
implementation at \texttt{sphglltools}); no closed form exists, unlike the
Chebyshev $\cos(j\pi/\modeidx_{\max})$ nodes, but this is a solved problem
requiring no new derivation. The associated LGL quadrature weights are
\begin{equation}
    w_j = \frac{2}{\modeidx_{\max}(\modeidx_{\max}+1)\,\big[P_{\modeidx_{\max}}(\yco_j)\big]^2} ,
    \label{eq:lgl-weights}
\end{equation}
with $w_j\cdot\measure(\yco_j,\Ny)$ the discrete quadrature weight for any
integral against the physical measure $\measure(\yco,\Ny)\,\d\yco$ of
Section~\ref{sec:msr-action} --- in particular for the endpoint,
$w_{\modeidx_{\max}} = 2/[\modeidx_{\max}(\modeidx_{\max}+1)]$, using
$P_{\modeidx_{\max}}(1)=1$: a known, finite number, needed below for the
terminal condition. Standard LGL differentiation matrices $D,D^{(2)}$
(first and second derivative on the nodes) are built once from the grid
alone --- they do not depend on $\Deltas(\Ny)$ or any other physics input,
so are computed once at the start and reused at every $\Ny$.

\subsection{Discretizing $\Lop$}

$\Lop$ (\eqref{eq:Lop-definition}) has $\Ny$-dependent, but not
$\Ny$-dependent-\emph{basis}, coefficients: at each $\Ny$, discretizing it
is
\begin{equation}
    \big(\Lop\field\big)_j
    =
    \e{\Deltas(\Ny)}\e{\Deltas(\Ny)\yco_j}
    \left[\frac{4}{\Deltas(\Ny)^2}\big(D^{(2)}\field\big)_j - \frac{2}{\Deltas(\Ny)}\big(D\field\big)_j\right] ,
\end{equation}
matching \eqref{eq:Lop-definition} --- note $\Lop$ itself does \emph{not}
carry the $1/\rout^2$ prefactor; that is applied separately, together with
$1/\aHloc^2$, wherever $\Lop$ appears in the equations of motion (e.g.\
eq. 8.1b).
i.e.\ $D,D^{(2)}$ applied to the node-value vector, then multiplied by
node-dependent \emph{scalar} coefficients recomputed at each $\Ny$ --- the
same pattern already used for every other solution-dependent coefficient
in this scheme ($\Hub^2\atloc$, $\epsSR\atloc$, $\Dcm{ij}\atloc$), no new
numerical primitive. The advection term $\advcoef(\yco,\Ny)\partial_\yco(\cdot)$
is likewise just $\mathrm{diag}(\advcoef(\yco_j,\Ny))\,D$ applied to the
node-value vector.

\subsection{Boundary conditions}

At $\yco=-1$ ($j=0$): ordinary Dirichlet rows for $\field,\mom$ (fixed to
$\phinl(\Ny),\pinl(\Ny)$) and $\rfield,\rmom$ (fixed to $0$) --- trivial.

At $\yco=+1$ ($j=\modeidx_{\max}$): Neumann rows for $\field,\rmom$
($\partial_\yco\field=0$, $\partial_\yco\rmom=0$), via \emph{hard
elimination}: solve $\sum_k D_{\modeidx_{\max},k}\,\field_k=0$ for the
boundary value itself,
\begin{equation}
    \field_{\modeidx_{\max}} = -\frac{1}{D_{\modeidx_{\max},\modeidx_{\max}}}\sum_{k\neq\modeidx_{\max}} D_{\modeidx_{\max},k}\,\field_k ,
\end{equation}
and substitute wherever $\field_{\modeidx_{\max}}$ appears in the interior
rows of $D^{(2)}$ (dense, so it appears in all of them). $\mom(1,\Ny)$ and
$\rfield(1,\Ny)$ are left unconstrained at this node --- they are read off
/ determined by the shooting problem (Section~\ref{sec:shooting}), not
fixed by a boundary row.

\begin{panel}
\semibold{Numerical scheme: hard elimination first.} Plain elimination on
the untransformed $\field,\mom,\rfield,\rmom$ gives no stability
guarantee, but is the cheap thing to try first --- build it, check
empirically whether stiff, spurious, or complex eigenvalues actually show
up in practice (the generic risk with Legendre/Chebyshev-type second-order
differentiation matrices under explicit time-stepping). If it works,
stop, until a case is found where it fails. If it fails, keep it as a
cross-check and build the alternative below.
\end{panel}

\begin{panel}
\semibold{Alternative: weighted SBP with SAT boundary terms.} A discrete
operator $D$ is SBP with respect to a norm matrix $H$ if
$HD+(HD)^T=\mathrm{diag}(-1,0,\dots,0,1)$ --- the discrete mirror of
integration by parts. For LGL nodes with the \emph{flat} measure, $H$ is
simply $\mathrm{diag}(w_j)$ and the SBP property holds essentially for
free; this is the standard reason LGL is the default choice where
provable stability matters. Our operator, however, is self-adjoint with
respect to the \emph{weighted}, $\Ny$-dependent measure $\measure(\yco,\Ny)\,\d\yco$
(Section~\ref{sec:msr-action}), not the flat one --- so $H$ here would
need to be $\mathrm{diag}(w_j\,\measure(\yco_j,\Ny))$, itself
$\Ny$-dependent, and the SBP property would need re-verifying (or the
operator re-derived in weighted form) at every $\Ny$, not just its
coefficients updated. This is more work than the flat-measure case makes
it look, though still standard technique (weighted/generalized SBP
operators are documented in the literature) rather than a research
problem. Boundary conditions would then be imposed weakly via
simultaneous approximation terms (SAT): rather than eliminating a
boundary unknown, add a penalty term to its RHS with strength fixed by the
SBP energy method, so that the discrete construction inherits a stability
estimate mirroring the continuum one. Build this only if plain elimination
turns out to need it.
\end{panel}

% ======================================================================
\section{Collocation scheme}
\label{sec:collocation-scheme}
% ======================================================================

\subsection{Grid representation}

$\field,\mom,\rfield,\rmom$ are represented directly by their node values
$\field_j(\Ny)\equiv\field(\yco_j,\Ny)$, etc. --- no mode expansion, and no
lift against the noiseless trajectory is structurally required.
Instead, in
collocation, the Dirichlet boundary row simply sets $\field_0(\Ny)=\phinl(\Ny)$
(Section~\ref{sec:collocation-basis}) --- no vanishing-basis obstruction
to work around. Tracking the deviation $\field_j-\phinl$ rather than
$\field_j$ directly may still be worth doing for numerical conditioning
(staying close to a well-behaved reference trajectory rather than
absolute field values that can range over many orders of magnitude), but
this is an optional implementation choice, and is deferred to the companion planning document.

\subsection{Equations of motion}

Substituting the grid representation into \eqref{eq:instanton-eqs} gives,
at each interior node $j$,
\begin{subequations}
\label{eq:colloc-eqs}
\begin{align}
    \dot\field_j
    & =
    \mom_j + \big(\mathrm{diag}(\advcoef)D\field\big)_j + \Dnoise{\field}(\yco_j,\Ny)\,\rfield_j + \Dnoise{\field\mom}(\yco_j,\Ny)\,\rmom_j ,
    \\
    \dot\mom_j
    & =
    -\big(3-\epsSR(\yco_j,\Ny)\big)\mom_j - \frac{\Vpot'(\field_j)}{\Hub^2(\yco_j,\Ny)}
    + \frac{1}{\rout^2\aHloc(\yco_j,\Ny)^2}\big(\Lop\field\big)_j
    + \big(\mathrm{diag}(\advcoef)D\mom\big)_j
    + \Dnoise{\mom}(\yco_j,\Ny)\,\rmom_j + \Dnoise{\field\mom}(\yco_j,\Ny)\,\rfield_j ,
    \\
    \dot{\rfield}_j
    & =
    -\frac{\Vpot''(\field_j)}{\Hub^2(\yco_j,\Ny)}\rmom_j
    + \frac{1}{\rout^2\aHloc(\yco_j,\Ny)^2}\big(\Lop\rmom\big)_j
    - \big(\mathrm{diag}(\advcoef)D\rfield\big)_j
    - \rfield_j\,\big(1-\epscore(\Ny)\big)\left[\frac{1}{\Deltas(\Ny)}-\frac32(\yco_j+1)\right] ,
    \\
    \dot{\rmom}_j
    & =
    \rfield_j - \big(3-\epsSR(\yco_j,\Ny)\big)\rmom_j
    - \big(\mathrm{diag}(\advcoef)D\rmom\big)_j
    - \rmom_j\,\big(1-\epscore(\Ny)\big)\left[\frac{1}{\Deltas(\Ny)}-\frac32(\yco_j+1)\right] ,
\end{align}
\end{subequations}
where $H^2(\yco_j,\Ny)\equiv\Hub^2[\field_j,\mom_j]$,
$\epsSR(\yco_j,\Ny)\equiv\epsSR[\field_j,\mom_j]$ --- ordinary pointwise
evaluation at each node, no reconstruct/project round-trip needed: 
the grid values \emph{are} already point
values.

\subsection{Terminal condition on the boundary node}
\label{sec:terminal-collocation}

The Lagrange-multiplier construction of Section~\ref{sec:bcs} carries
over directly, and is strikingly simple.
Varying $S_{\mathrm{tot}}=S+\lag(\field_{\modeidx_{\max}}(\Nfinal)-\phiend)$
with respect to the node values $\{\field_j(\Nfinal)\}$ gives, using the
quadrature weights of \eqref{eq:lgl-weights},
\begin{equation}
    \sum_j w_j\,\measure(\yco_j,\Nfinal)\,\rfield_j(\Nfinal)\,\delta\field_j
    + \lag\,\delta\field_{\modeidx_{\max}}
    = 0
    \qquad \text{for all admissible } \{\delta\field_j\} .
\end{equation}
Since $\{\delta\field_j\}$ are already independent, ordinary
finite-dimensional variables,
this holds for every $j$ separately:
\begin{equation}
    \rfield_j(\Nfinal) = 0 \quad (j\neq\modeidx_{\max}) ,
    \qquad
    \rfield_{\modeidx_{\max}}(\Nfinal) = -\frac{\lag}{w_{\modeidx_{\max}}\,\measure(1,\Nfinal)} ,
    \label{eq:terminal-colloc}
\end{equation}
using $w_{\modeidx_{\max}}=2/[\modeidx_{\max}(\modeidx_{\max}+1)]$ from
\eqref{eq:lgl-weights}. Eq.~\eqref{eq:terminal-colloc} is
the exact, complete terminal condition, not an asymptotic or truncated
approximation to one.

\subsection{Shooting problem}
\label{sec:shooting}

Closed by a single scalar shooting condition: adjust $\lag$ so that
\begin{equation}
    \field_{\modeidx_{\max}}(\Nfinal) = \phiend .
\end{equation}
Structurally identical to the shooting problem already solved for the
single-trajectory instanton.

\subsection{Picard iteration}

\begin{enumerate}
    \item Forward pass: integrate $(\field_j,\mom_j)$ forward from
    \eqref{eq:bc-init} with $\rfield_j=\rmom_j=0$ (zeroth iterate).
    \item Backward pass: integrate $(\rfield_j,\rmom_j)$ backward from
    \eqref{eq:terminal-colloc}, using the current forward-pass solution to
    evaluate the nonlinear coupling $\propto\Vpot''(\field)$ and the
    local $\Hub^2,\epsSR,\Dcm{ij},\advcoef$ coefficients.
    \item Forward pass: re-integrate $(\field_j,\mom_j)$ using the
    updated response-field values as sources.
    \item Repeat to convergence; adjust $\lag$ via an outer scalar
    root-find to satisfy the shooting condition.
\end{enumerate}

% ======================================================================
\section{Open issues}
\label{sec:open-issues}
% ======================================================================

\subsection{Sub-horizon noise: resolved}

An earlier draft carried this as the first, blocking open issue: treating
every $\yco$ interior to the current horizon as an independently
$\delta(\yco-\yco')$-correlated noise source manufactures degrees of
freedom that do not physically exist, since $\ncount(\yco,\Ny)\to0$ there
(fewer than one coarse-graining patch per shell). Three options were
weighed --- switching off noise there (fails the reduction-limit
cross-check), enforcing a single shared noise realization there (a hard,
time-varying shape constraint on the solution), and restricting $\yco$ to
the exterior of the horizon (previously rejected due to a
time-dependent eigenbasis) --- and a rank-1 shared-noise-channel
modification of the MSR action was proposed as a leading candidate,
explicitly flagged as unpropagated through the full equation set.

That proposal is now retired. Adopting the exterior-only, logarithmic
onion coordinate of Section~\ref{sec:coordinate}, combined with a
collocation (rather than explicit spectral) numerical method, removes the
objection to horizon-exclusion that motivated rejecting it originally,
and the problem is resolved \emph{by construction} rather than patched:
there is no sub-horizon region left in the domain to source spurious
independent fluctuations, so no special noise treatment is needed at all.
See Section~\ref{sec:coordinate} for the coordinate definition and the
panel explicitly narrating this reversal.

\subsection{Probability interpretation of $S_{\mathrm{MSR}}$ for the 2D configuration}

\begin{todo}
The single-trajectory instanton action $S_{\mathrm{MSR}}^{\mathrm{1D}}$
computes the probability of a specific event: the central Hubble patch
follows a given trajectory, \emph{with the surrounding shells assumed to
evolve noiselessly} (zero action cost by assumption). The two-dimensional
action $S_{\mathrm{MSR}}$ of \eqref{eq:msr-action} computes the probability
of the joint configuration $\field(\yco,\Ny)$ across the whole overdensity,
core and shells together. These are not simply refinements of the same
number: they are probabilities for two different events. Naively, one
might attempt to interpret the difference $S_{\mathrm{MSR}} -
S_{\mathrm{MSR}}^{\mathrm{1D}}$ as `the additional cost of the gradient
correction', but it was pointed out (correctly) in discussion that the
outer shells are not passive bystanders: they participate directly in the
compaction function and hence in the collapse criterion, exactly as they
do in the existing (noiseless-detachment) construction. So it is not
obviously right to think of the 1D action as already including a
`for-free' contribution from the outer layers that the 2D calculation
merely refines --- rather, the 1D action already implicitly describes the
joint configuration too, just with the outer-layer part of that
configuration fixed to be exactly noiseless. The correct comparison, and
the correct interpretation of what $S_{\mathrm{MSR}}$ computes once the
gradient coupling is switched on, has not yet been fully resolved and
needs further thought before results can be presented as `the corrected
PBH formation probability'.
\end{todo}

\subsection{The largest time equation for a stochastic (MSR) path integral}

\begin{todo}
The terminal condition $\rfield(\yco,\Nfinal)=0$ (Section~\ref{sec:bcs})
was justified heuristically by an analogue of the CTP `largest time
equation': placing the true response-field terminal condition at some
$\Ny_{\mathrm{turn}} > \Nfinal$ should not change the solution for
$\Ny \leq \Nfinal$, provided the interval $[\Nfinal,\Ny_{\mathrm{turn}}]$ is
genuinely noiseless and gradient-free, because the response-field
equations of motion are then linear and homogeneous there and, starting
from zero at $\Ny_{\mathrm{turn}}$, remain zero all the way back to
$\Nfinal$. This is physically reasonable and was checked at the level of
the linearized equations here, but the underlying general statement ---
that a Schwinger--Keldysh-type path integral, cut open at an intermediate
time with no later operator insertions, collapses to a functional delta
function enforcing agreement of the forward and backward fields at that
time --- does not appear to have a fully satisfactory published derivation
even in the standard quantum CTP case (see the parallel discussion in the
companion notes on the Dyson equation and the $\delta N$/Schwinger--Keldysh
correspondence), and has certainly not been established here for the
\emph{stochastic} (MSR) path integral in the presence of nonlinear noise
and gradient coupling. The quantum-case argument relies on unitarity of
the $[\Nfinal,\Ny_{\mathrm{turn}}]$ evolution; the analogous statement in
the stochastic case would rest on probability conservation, but the
precise functional-integral statement of this, and whether it implies the
same clean factorization, has not been derived carefully here. A
worthwhile check, cheaper than a full derivation: solve the existing
single-trajectory (1D) instanton with the terminal condition placed at
several $\Ny_{\mathrm{turn}} > \Nfinal$ and verify numerically that the
solution for $\Ny \leq \Nfinal$ is insensitive to the choice, before
relying on the same argument in the 2D case.
\end{todo}

% ======================================================================
\section{Summary}
\label{sec:summary}
% ======================================================================

The gradient-corrected instanton is the saddle point of the action
\eqref{eq:msr-action} on the fixed rectangle $\yco\in[-1,1]$
(exterior to the current horizon only), $\Ny\in[\Ninit,\Nfinal]$,
equivalently the collocation system \eqref{eq:colloc-eqs} with boundary
conditions \eqref{eq:bc-init}--\eqref{eq:terminal-colloc} and shooting
parameter $\lag$, solved by Picard iteration
(Section~\ref{sec:collocation-scheme}). The corrected core trajectory
$\field(1,\Ny)$ and the full profile $\zprofile(\yco)$ of
\eqref{eq:zeta-extraction} are outputs of the solution, to be compared
against the existing noiseless-detachment construction. Two theoretical
issues (Section~\ref{sec:open-issues}) remain open and should be resolved,
or at least understood well enough to bound their impact, before
quantitative results are reported: the probability interpretation of the
2D action, and the largest-time-equation justification of the terminal
boundary condition. The previously-blocking sub-horizon-noise issue is
resolved by the exterior-only coordinate itself, not by a separate
numerical patch. A discrete (finite-shell) implementation
(Section~\ref{sec:geometry}) remains available as an independent numerical
cross-check on the collocation scheme of
Sections~\ref{sec:collocation-basis}--\ref{sec:collocation-scheme}.

\end{document}
